\section*{Hoofdstuk \ref{ch:g12:polymers} --- Polymeren}

\begin{solution}{exo:g12:polymers:1}
Propeen, \ce{CH2=CH-CH3}.
\end{solution}

\begin{solution}{exo:g12:polymers:2}
\ce{-CF2-CF2-}; een \omterm{def:g12:polymers:addition-polymer}{additiepolymeer}: het \omterm{def:g12:polymers:polymer}{monomeer} heeft een
\ce{C=C}-binding en er ontstaat geen ander product.
\end{solution}

\begin{solution}{exo:g12:polymers:3}
Additie: polyetheen, PVC, polystyreen. Condensatie: PET (een polyester),
nylon (een polyamide).
\end{solution}

\begin{solution}{exo:g12:polymers:4}
Cellulose, zetmeel, rubber uit latex, eiwitten. Nylon en PVC zijn
synthetisch.
\end{solution}

\begin{solution}{exo:g12:polymers:5}
Elke koppeling gebruikt één groep van elk \omterm{def:g12:polymers:polymer}{monomeer}; met twee groepen kan
een \omterm{def:g12:polymers:polymer}{monomeer} aan beide kanten binden en blijft de keten groeien. Een
\omterm{def:g12:polymers:polymer}{monomeer} met één groep zou de keten afsluiten.
\end{solution}

\begin{solution}{exo:g12:polymers:6}
$M(\ce{C2H3Cl}) = 24.0 + 3.0 + 35.5 = \qty{62.5}{g/mol}$;
$n = 125000/62.5 = 2000$.
\end{solution}

\begin{solution}{exo:g12:polymers:7}
$M(\ce{C8H8}) = \qty{104.0}{g/mol}$; $M = 2500 \times 104.0 =
\qty{2.60e5}{g/mol}$.
\end{solution}

\begin{solution}{exo:g12:polymers:8}
De zuurgroep van het ene \omterm{def:g7:atoms-and-molecules:molecule}{molecuul} vormt een \omterm{def:g11:functional-groups:carboxylic-acid}{ester} met de alcoholgroep van
het volgende, dat nog een vrije zuurgroep heeft: de keten groeit.
\omterm{def:g12:polymers:repeat-unit}{Repeterende eenheid} \ce{-O-CH(CH3)-CO-}; bij elke koppeling wordt water
afgesplitst.
\end{solution}

\begin{solution}{exo:g12:polymers:9}
Folie voor tassen is gemaakt van vertakte ketens, die niet kunnen
stapelen: het materiaal is amorf, zacht en doorzichtig. Melkflessen zijn
gemaakt van bijna lineaire ketens, die stapelen tot kristallijne
gebieden: dichter, stijver, en de kristallijne gebieden verstrooien
licht.
\end{solution}

\begin{solution}{exo:g12:polymers:10}
Hun ketens zijn door veel covalente dwarsverbindingen aan elkaar
vastgemaakt. Verwarmen kan de ketens niet langs elkaar laten glijden; het
verbreekt bindingen en vernietigt het materiaal in plaats van het te
smelten.
\end{solution}

\begin{solution}{exo:g12:polymers:11}
$M(\ce{C12H22N2O2}) = 144.0 + 22.0 + 28.0 + 32.0 = \qty{226.0}{g/mol}$;
$n = 22600/226.0 = 100$.
\end{solution}

\begin{solution}{exo:g12:polymers:12}
\[
  \ce{H2N(CH2)6NH2 + HOOC(CH2)4COOH -> H2N(CH2)6NHCO(CH2)4COOH + H2O} .
\]
Het product heeft nog een vrije aminogroep aan het ene uiteinde en een
vrije zuurgroep aan het andere: elk kan met een ander \omterm{def:g12:polymers:polymer}{monomeer} reageren.
\end{solution}

\begin{solution}{exo:g12:polymers:13}
Twee \omterm{def:g7:atoms-and-molecules:molecule}{moleculen} water per \omterm{def:g12:polymers:repeat-unit}{repeterende eenheid}. $n = 1000/226.0 =
\qty{4.42}{mol}$ \omterm{def:g12:polymers:repeat-unit}{repeterende eenheden}, dus \qty{8.85}{mol} water:
$8.85 \times 18.0 = \qty{159}{g}$.
\end{solution}

\begin{solution}{exo:g12:polymers:14}
Warm natuurrubber bestaat uit ketens die langs elkaar glijden.
Zwavelbruggen verbinden de ketens: ze rekken uit maar veren terug, en
vloeien niet meer. Omdat een band dwarsverbindingen heeft, kun je hem niet
smelten en opnieuw vormen.
\end{solution}

\begin{solution}{exo:g12:polymers:15}
$M(\ce{C6H10O5}) = \qty{162.0}{g/mol}$; $n = \num{1.6e6}/162.0 = 9900$
eenheden; lengte $9900 \times 0.5 = \qty{4900}{nm}$, ongeveer
\qty{5}{\micro m}.
\end{solution}

\begin{solution}{pb:g12:polymers:1}
\textbf{1.} Ethaan-1,2-diol (twee alcoholgroepen) en
benzeen-1,4-dicarbonzuur (twee carbonzuurgroepen).

\textbf{2.} Een esterbinding; er wordt water afgesplitst.

\textbf{3.} Een \omterm{def:g12:polymers:addition-polymer}{condensatiepolymeer}.

\textbf{4.} $M(\ce{C2H6O2}) = \qty{62.0}{g/mol}$; $M(\ce{C8H6O4}) =
\qty{166.0}{g/mol}$.

\textbf{5.} $62.0 + 166.0 - 2 \times 18.0 = \qty{192.0}{g/mol}$; en
$10 \times 12.0 + 8 \times 1.0 + 4 \times 16.0 = 192.0$.

\textbf{6.} $n = 38400/192.0 = 200$.

\textbf{7.} Twee esterbindingen per \omterm{def:g12:polymers:repeat-unit}{repeterende eenheid}: ongeveer 400.

\textbf{8.} Zijn ketens zijn grotendeels lineair, en delen ervan liggen
op een rij in geordende gebieden.

\textbf{9.} Het is een \omterm{def:g8:plastics:thermoplastic}{thermoplast}: zijn ketens hebben geen
dwarsverbindingen en glijden langs elkaar als ze heet zijn.

\textbf{10.} $300/0.80 = \qty{375}{g}$.

\textbf{11.} $375/25 = 15$ flessen.

\textbf{12.} Doppen, etiketten, vuil en snippers worden verwijderd of
gaan verloren tijdens het sorteren, wassen en spinnen.

\textbf{13.} Voorkom afval (principe 1): een gebruikt voorwerp wordt
\omterm{def:g5:raw-materials-and-recycling:raw-material}{grondstof}.

\textbf{14.} \ce{C10H8O4 + 2H2O -> C2H6O2 + C8H6O4}.

\textbf{15.} $1000/192.0 = \qty{5.21}{mol}$.

\textbf{16.} Diol: $5.21 \times 62.0 = \qty{323}{g}$; zuur:
$5.21 \times 166.0 = \qty{865}{g}$.

\textbf{17.} $2 \times 5.21 \times 18.0 = \qty{188}{g}$ water;
$1000 + 188 = 323 + 865 = \qty{1188}{g}$: de massa blijft behouden.

\textbf{18.} De \omterm{def:g12:polymers:polymer}{monomeren} kunnen gezuiverd worden en geven nieuw PET dat
zo goed is als nieuw, terwijl smelten de ketens telkens een beetje
korter maakt.

\textbf{19.} Voor één fleecevest zijn 15 flessen van \qty{25}{g} nodig.
\end{solution}
